\documentclass[10pt]{article}
\usepackage[margin=1in]{geometry}
\usepackage{amsmath,amssymb,amsthm,mathtools}
\usepackage[hidelinks]{hyperref}
\newtheorem{theorem}{Theorem}
\newtheorem{lemma}[theorem]{Lemma}
\newcommand{\Z}{\mathbb Z}
\newcommand{\Q}{\mathbb Q}
\title{Denominators of two basic hypergeometric ratios\\
\large A proof of MathDB problem 346880}
\author{Principia Math In Progress}
\date{September 12, 2026}
\begin{document}
\maketitle
\noindent\textbf{Task 1: solve the original problem.}
We prove Conjecture 8.4 of Kim--Stanton \cite{KS}, which is the
MathDB question, and its companion Conjecture 8.3. These are
Conjectures 7.3 and 7.2, respectively, in the earlier manuscript
\cite{original}. Only one MathDB problem is counted.

Let $a,b,c,q$ be independent indeterminates and write
\[
 (u;q)_j=\prod_{i=0}^{j-1}(1-uq^i),\qquad
 F(a,b,c;z)={}_2\phi_1(a,b;c;q,z)
   =\sum_{j\ge0}\frac{(a;q)_j(b;q)_j}{(c;q)_j(q;q)_j}z^j.
\]
Define formal series
\[
 R(z)=\frac{F(aq,b,cq;z)}{F(a,b,c;z)}
       =\sum_{n\ge0}r_nz^n,\qquad
 S(z)=\frac{F(aq,bq,cq;z)}{F(a,b,c;z)}
       =\sum_{n\ge0}s_nz^n .
\]
Their constant coefficients are $r_0=s_0=1$. Put
\[
 D_n(c,q)=\prod_{k=0}^n(1-cq^k)^{\lfloor(n+1)/(k+1)\rfloor}.
\]
\begin{theorem}\label{main}
For every integer $n\ge0$,
\[
 \frac{D_n(c,q)}{1-c}\,r_n\in\Z[a,b,c,q],
 \qquad
 \frac{D_n(c,q)}{1-c}\,s_n\in\Z[a,b,c,q].
                                                               \tag{1}
\]
The first assertion is the second coefficient-denominator conjecture
recorded as MathDB 346880.
\end{theorem}

We first remove possible denominators involving $q$ alone. Set
\[
 B=\Z[a,b,q],\quad A=B[c],\quad
 \mathcal L=A[(1-cq^j)^{-1}:j\ge0],\quad K=\Q(a,b,q).
\]
\begin{lemma}\label{integral}
Every coefficient of $R$ and $S$ belongs to $\mathcal L$.
Moreover,
\[
 \mathcal L\cap K[c]=A.                                  \tag{2}
\]
\end{lemma}
\begin{proof}
Heine's contiguous relation, also given in \cite[p.~13]{KS}, is
\[
 F(aq,b,cq;z)-F(a,b,c;z)
   =\beta z F(aq,bq,cq^2;z),\qquad
 \beta=\frac{(a-c)(1-b)}{(1-c)(1-cq)}.                    \tag{3}
\]
It can be checked directly in each coefficient: the relevant numerator
simplification is
\[
 (1-au)(1-c)-(1-a)(1-cu)=(a-c)(1-u),\qquad u=q^j,
\]
and $1-q^j$ cancels from $(q;q)_j$.
Let $\sigma$ substitute $(a,b,c)\mapsto(b,aq,cq)$, leaving $q,z$
fixed. Symmetry of $F$ in its first two parameters changes (3) into
\[
 R=\frac1{1-\beta z\,\sigma(R)},\qquad
 r_n=\beta\sum_{j=0}^{n-1}r_j\,\sigma(r_{n-1-j})\quad(n\ge1).
                                                               \tag{4}
\]
Both $\beta\in\mathcal L$ and
$\sigma(\mathcal L)\subseteq\mathcal L$ hold.
Induction on $n$ proves $r_n\in\mathcal L$.

A second elementary contiguous identity is
\[
 F(aq,bq,cq;z)-F(aq,b,cq;z)
   =\delta z F(aq^2,bq,cq^2;z),\qquad
 \delta=\frac{b(1-aq)}{1-cq}.                             \tag{5}
\]
For a direct verification, in the coefficient of $z^j$ the shift
$b\mapsto bq$ contributes $b(1-q^j)/(1-b)$; cancellation in the
shifted factorials gives the right side of (5).
If $\tau$ substitutes $(a,b,c)\mapsto(aq,b,cq)$, then
\[
 S=R\bigl(1+\delta z\,\tau(S)\bigr),\qquad
 s_n=r_n+\delta\sum_{j=0}^{n-1}r_j\,\tau(s_{n-1-j}).
                                                               \tag{6}
\]
Since $\tau$ also preserves $\mathcal L$, induction proves
$s_n\in\mathcal L$.

Finally, write an element of $\mathcal L$ as $U/V$ with $U\in A$
and $V$ a finite product of powers of $1-cq^j$. Its constant
coefficient in $c$ is $V(0)=1$, so its Taylor series at $c=0$ has
all coefficients in $B$. If $U/V$ is also a polynomial in $K[c]$,
that Taylor series is finite. The polynomial therefore lies in
$B[c]=A$, proving (2).
\end{proof}

\begin{proof}[Proof of Theorem \ref{main}]
Let $H$ denote either $F(aq,b,cq;z)$ or $F(aq,bq,cq;z)$,
and put $F=F(a,b,c;z)$. We bound the poles, as functions of $c$,
of the coefficients of
\[
 \frac{H}{(1-c)F}.                                      \tag{7}
\]
Work over $K=\Q(a,b,q)$. The points $c=q^{-k}$, $k\ge0$, are
distinct, and each factor $1-cq^k$ is linear in $c$.

Fix $k\ge1$ and put $\epsilon=1-cq^k$.
The coefficient of $z^j$ in $F$ is regular at $c=q^{-k}$ for
$j\le k$, and has at most a simple pole for $j\ge k+1$.
For $H$ the corresponding threshold is $j=k$: its coefficients
are regular for $j<k$ and have at most a simple pole thereafter.
Thus, over the local ring of rational functions regular at this point,
\[
 \epsilon F=\epsilon U+z^{k+1}V,\qquad
 \epsilon H=\epsilon W+z^kX,                             \tag{8}
\]
where $U,W,V,X$ are formal series with regular coefficients,
and $U(0)=W(0)=1$. In particular $U$ is a unit in the formal
series ring over the local ring. Expand
\[
 \frac HF=
 \left(W+\frac{z^kX}{\epsilon}\right)\frac1U
 \sum_{\ell\ge0}
       \frac{(-1)^\ell z^{(k+1)\ell}V^\ell}
                            {\epsilon^\ell U^\ell}.     \tag{9}
\]
For each coefficient of $z^n$ only finitely many summands contribute.
The part involving $W$ has pole order at most
$\lfloor n/(k+1)\rfloor$. The other part vanishes if $n<k$;
otherwise its pole order is at most
\[
 1+\left\lfloor\frac{n-k}{k+1}\right\rfloor
   =\left\lfloor\frac{n+1}{k+1}\right\rfloor.             \tag{10}
\]
The factor $1-c$ is a local unit because $k\ge1$, so (10) is
also the required bound for (7). In particular there is no pole
when $k>n$.

The point $c=1$ requires separate treatment. Now
$\epsilon=1-c$, all coefficients of $H$ are regular, and
\[
 \epsilon F=\epsilon+zV
\]
with regular coefficients in $V$. Hence (7) becomes
\[
 \frac{H}{\epsilon+zV}
   =H\sum_{\ell\ge0}\frac{(-1)^\ell z^\ell V^\ell}
                           {\epsilon^{\ell+1}}.         \tag{11}
\]
Its coefficient of $z^n$ has pole order at most $n+1$,
which is precisely the exponent prescribed at $k=0$.

Lemma \ref{integral} shows that the coefficients of (7) lie in
$\mathcal L$, so there are no other possible finite poles in $c$.
Equations (10)--(11) show that $D_n$ removes every such pole.
The resulting rational function therefore lies in $K[c]$.
It also lies in $\mathcal L$, so (2) places it in
$A=\Z[a,b,c,q]$, as required.
\end{proof}

For example, the normalized polynomial for $R$ at $n=1$ is
$(a-c)(1-b)$, and both normalized polynomials at $n=0$ are $1$.
The proof is an identity over the field of rational functions;
specializations are understood wherever the original series coefficients
are defined, or through the resulting rational identities.

\medskip
\noindent\textbf{Task 2: determine the minimal denominators.}
Compute the exact generic pole orders of $r_n/(1-c)$ and $s_n/(1-c)$
at every $c=q^{-k}$. Determine when the exponents in $D_n$ are attained
and classify parameter specializations that lower them.

\medskip
\noindent\textbf{Task 3: treat arbitrary integral parameter shifts.}
For nonnegative integers $u,v,w$, study
\[
 F(aq^u,bq^v,cq^w;z)/F(a,b,c;z).
\]
Determine which shifts admit
normalizations in $\Z[a,b,c,q]$ and derive their optimal pole bounds.

\medskip
\noindent\textbf{Task 4: give explicit formulas for the normalized numerators.}
Construct finite sums or combinatorial models for
$D_nr_n/(1-c)$ and $D_ns_n/(1-c)$ in which denominator cancellation
has already been carried out. Use them to explain systematic factors
and coefficient patterns.

\medskip
\noindent\textbf{Task 5: extend the method to higher basic hypergeometric series.}
Identify contiguous ratios of \({}_{r+1}\phi_r\) whose coefficients lie
in an integral localization of their lower parameters, and establish
denominator theorems by combining that property with local pole bounds.

\begin{thebibliography}{9}
\bibitem{KS}
J.~S.~Kim and D.~Stanton,
\emph{Three families of $q$-Lommel polynomials},
S\'eminaire Lotharingien de Combinatoire \textbf{87B} (2023),
Article 8, 25 pp.; Conjectures 8.3--8.4.
\href{https://www.mat.univie.ac.at/~slc/wpapers/lattpath/KimStanton/KimStanton.pdf}
{Published full text}.
\bibitem{original}
J.~S.~Kim and D.~Stanton,
\emph{Ratios of Hahn--Exton $q$-Bessel functions and $q$-Lommel
polynomials}, arXiv:2006.08120v2 (2021), Conjectures 7.2--7.3.
\href{https://arxiv.org/abs/2006.08120}{arXiv record}.
\end{thebibliography}
\end{document}
